\documentclass[a4paper,addpoints,14pt]{exam}
%\documentclass[answers,addpoints,14pt]{exam}
\unframedsolutions

\usepackage{array,amsmath,fancybox,amsfonts,graphicx,amsthm,amssymb,wasysym}
\usepackage{latexsym,slashed,anysize,subfigure,calc}
\usepackage[spanish]{babel}
\usepackage[utf8]{inputenc}
%\usepackage{fancyhdr,float,hyperref}
%\usepackage{lettrine}

%\usepackage{enumitem}



%para los entornos de soluciones
\usepackage{mdframed,color}
	%uso: 	entorno (begin,end): {mdframed}[backgroundcolor=brown!10] 



%entorno para química: texto y math mode
%\newcommand*\chem[1]{\textup{#1}}
\newcommand*\chem[1]{\ensuremath{\mathrm{#1}}}

%tamaños fuentes extra: 8,9,14,17,20
	\usepackage{extsizes}

%display style in long fractions
\newcommand\ddfrac[2]{\frac{\displaystyle #1}{\displaystyle #2}}

%tabla generada
\usepackage[table,xcdraw]{xcolor}

%lewis y química
\usepackage{chemformula,elements}
%chemformula, provides \ch{} with the !(<below>)(<formula>) syntax and \chlewis;
%elements, provides \elconf and \writeelconf.
%ex:
%\ch{
%  !(\elconf{Li})( "\chlewis{180.}{Li}" ) +
%  !(\elconf{F})( "\chlewis{0.90:180:270:}{F}" )
%   ->
%  !(\writeelconf{2})( Li+ ) +
%  !(\writeelconf{2,2+6})( "\chlewis{0:90:180:270:}{F}" {}- )
%}

%multicol
\usepackage{multicol}
  

%símbolo bonito para qed
\usepackage{amsthm}
	%\newcommand{\mc}[3]{\multicolumn{#1}{#2}{#3}}
\newcommand{\halmos}{\vspace*{-6mm}
	
					\begin{flushright}
						\rule{1mm}{2.5mm} 
					\end{flushright}
					\vspace*{-4mm}
						
					  }
\renewcommand{\qed}{\halmos}

 
%página blanco
\usepackage{afterpage}
%uso \afterpage{\blankpage}

\newcommand\blankpage{%
    \null
    \thispagestyle{empty}%
    \addtocounter{page}{0}%
    \newpage}





 %caja bonita
  \usepackage[many]{tcolorbox}
\definecolor{titlebg}{RGB}{80,100,72}
\newtcolorbox{learn}{
  breakable,
  enhanced,
  arc=0pt,
  outer arc=0pt,
  colframe=titlebg,
  colback=titlebg!7,
  overlay unbroken and first={
    \node[
      draw=titlebg,
      fill=titlebg,
      rotate=270,
      anchor=north west,
      text=white,
      font=\bfseries
    ]
    at (frame.north west)  
    {DATOS};
  }
}


\newtcolorbox{formulacion}{
  breakable,
  enhanced,
  arc=0pt,
  outer arc=0pt,
  colframe=titlebg,
  colback=titlebg!7,
  overlay unbroken and first={
    \node[
      draw=titlebg,
      fill=titlebg,
      rotate=270,
      anchor=north west,
      text=white,
      font=
    ]
    at (frame.north west)  
    {\hspace*{3ex}};
  }
}

\newtcolorbox{problemas}{
  breakable,
  enhanced,
  arc=0pt,
  outer arc=0pt,
  colframe=titlebg,
  colback=titlebg!7,
  overlay unbroken and first={
    \node[
      draw=titlebg,
      fill=titlebg,
      rotate=270,
      anchor=north west,
      text=white,
      font=
    ]
    at (frame.north west)  
    {AVISO};
  }
}
	
\usepackage{anysize}
	%\marginsize{left}{right}{top}{bottom}
    \marginsize{1cm}{.8cm}{1cm}{1cm}
    
\usepackage[spanish]{babel}
\usepackage[utf8]{inputenc}

\usepackage{sidecap}
\usepackage{mwe,float,afterpage}

%redefiniciones
\pointpoints{punto}{puntos}
\bonuspointpoints{punto}{puntos}
\hqword{Pregunta} %substitutes text for \Question"
%\vpgword{Pág.}%substitutestext for \Page"
\hpword{Máx.}%substitutes text for \Points"
\hsword{Nota}%substitutes text for \Score"
\htword{Total}%substitutes text for \Total:"

\chqword{Pregunta}
\chpword{Máx.}
\chsword{Nota}
\chtword{Total}
\chbpword{Puntos \textit{bonus}}

\renewcommand{\solutiontitle}{\noindent\textbf{Solución:}\par\noindent}


%fuentes
	%\usepackage[T1]{fontenc}
	%\usepackage{fourier}

	\usepackage[T1]{fontenc}
	\usepackage{librebaskerville}


\usepackage{graphicx,graphics}
\usepackage{color}
	%	\colorgrids %grids con color
	%	\definecolor{GridColor}{gray}{0.8}
		%grid
%			\setlength{\gridsize}{5mm}
%   			 \setlength{\gridlinewidth}{0.1pt}

%fracciones en texto más elaboradas
	\usepackage{xfrac}

%mejora en la exposición de guarismos
	\usepackage{numprint}
 
%caption
\usepackage[font=small,labelfont=bf,
labelsep=endash]{caption}
	
%soluciones sombreadas	
\shadedsolutions


\checkboxchar{$\Box$}
\checkedchar{$\blacksquare$}

\usepackage{everypage,draftwatermark} %Marca de agua
\SetWatermarkText{\hspace*{-4cm} \parbox{\textwidth}{\centering  B1\\dist.}}
\SetWatermarkScale{2}
\SetWatermarkLightness{0.9}
%
%\usepackage{background}
%\backgroundsetup{
%scale=1.5,
%angle=180,
%opacity=.05,
%color=black,
%contents={
%\includegraphics[scale=1]{unicorn.jpg}}}


%\pagestyle{headandfoot}
%\runningheadrule
%\firstpageheader{Física para Biólogos}{Boletín 1a}{1 de octubre de 2011}
%\runningheader{Física para Biólogos}
%              {Boletín 1a, Página \thepage\ de \numpages}
%            {1 de octubre de 2011}
%\firstpagefooter{}{}{}
%\runningfooter{}{}{}


%espacios extra
%\extrawidth{1in}
%%\extrafootheight{.75in}
%\extraheadheight{1in}

%cabecera
\header{\includegraphics[scale=.22]{a.png}}       
      {}
       {$f_2$\textbf{B1-{\small Distancia}}\\\textbf{febrero}2023}       
%pie de página

   %enumera en todas menos en la primera
   %\runningfootrule

   %enumera en todas
   %\footrule
   %\footer{}{Página \thepage\ de \numpages}{}
%   \footer{}
%	{Página \thepage\ de \numpages}
%	{\iflastpage{{\footnotesize \textit{Fin del examen}}}{{\footnotesize \textit{Da la vuelta a la hoja}\ldots}}}


		
\begin{document}
%portada
%formato de preguntas
\qformat{\textbf{Pregunta \thequestion}\quad (\thepoints)\hfill} 
\pagenumbering{gobble}

\fbox{\parbox{0.945\linewidth}{\textbf{Nombre y apellidos}:\vspace*{20mm}}}
%\vspace{1ex}

%\textbf{\hspace*{5mm} Sólo una de las respuestas es correcta.\newline \hspace*{11mm} Las preguntas erróneamente contestadas restarán 0.25 puntos.}
\vspace*{\fill}

\begin{learn}
{$\triangleright\; M(\chem{O})\simeq16\;g\,mol^{-1}$\newline
$\triangleright\; M(\chem{N})\simeq14\;g\,mol^{-1},$\newline
$\triangleright\; M(\chem{H})\simeq1\;g\,mol^{-1},$\newline
$\triangleright\; M(\chem{S})\simeq32\;g\,mol^{-1},$\newline
$\triangleright\; M(\chem{Cu})\simeq64\;g\,mol^{-1},$\newline
$\triangleright\; R\simeq 0.082\;atm\,L\,mol^{-1}\,K^{-1}$}
\end{learn}	
%		\vspace{0ex}
		

%\begin{question}[1] ¿Cuál de las siguientes afirmaciones es \textbf{\small{VERDADERA}}?
%	\begin{checkboxes} 
%			\choice El enlace covalente sólo se da entre metales.			
%			\choice Todos los compuestos covalentes se encuentran en estado sólido a temperatura ambiente.	
%			\CorrectChoice Los compuestos covalentes pueden formar redes cristalinas con puntos de ebullición muy altos.
%			\choice Los sólidos covalentes conducen la electricidad disueltos.	
%	\end{checkboxes}
%\end{question}
%\qed
%%
\newpage
\pagenumbering{gobble}

\begin{problemas}{Recuerda que un examen es una \textit{comunicación}: debes invocar las leyes en que te basas para resolver los ejercicios así como \textit{indicar brevemente los pasos y razonamientos} que sigues para obtener los resultados.}
\end{problemas}
\vspace{1ex}


\begin{questions}

\question[2] En un matraz de $1.0\;L$ se introducen $0.02\;mol$ de oxígeno molecular y $0.03\;mol$ de nitrógeno molecular. Si la temperatura es de $20^\circ C$ y asumimos que ambos gases se comportan idealmente, calcula \textit{razonadamente}:
	\begin{itemize}
		\item[a)] La presión total de la mezcla.
		
		\begin{mdframed}[backgroundcolor=brown!10]  
		La presión total puede hallarse calculando el número total de moles de la mezcla e imponiendo que cumplan la ecuación de estado
		\[P_TV=n_TRT\rightarrow P_T=\frac{(0.02+0.03)\cdot 0.082\cdot 293}{1}\simeq 1.2\;atm.\]
		\end{mdframed}
		
		\item[b)] La fracción molar y la presión parcial de cada gas.
		\begin{mdframed}[backgroundcolor=brown!10]  
		Las fracciones molares serán
		\begin{eqnarray*}
				\chi_{\chem{O_2}}&=&\frac{0.02}{0.02+0.03}\simeq 0.4,\\
				\chi_{\chem{N_2}}&=&\frac{0.03}{0.02+0.03}\simeq 0.6.
		\end{eqnarray*}
		y que conduce, en virtud de la ley de Dalton, a las presiones parciales
		\[P_{\chem{O_2}}\simeq 0.5\;atm;\hspace{1cm}P_{\chem{O_2}}\simeq 0.7\;atm.\] 
		\end{mdframed}
	\end{itemize}
\qed

\newpage
\pagenumbering{gobble}


\question[2] Estima la molaridad de:
	\begin{itemize}
		\item[a)]Una disolución de ácido nítrico etiqueda como al $16\%$ en peso y densidad $1.09\;g\,cm^{-3}$.
		\begin{mdframed}[backgroundcolor=brown!10]  
		Supongamos $1\;L$ de disolución: su masa es de $1090\;g$ y de ella, $174.4\;g$ ($2.8\;mol$) corresponden a ácido nítrico puro. Por tanto $$M\simeq \frac{2.8}{1}=2.8\;mol\,L^{-1}.$$
		
		\end{mdframed}
		\item[b)] La disolución resultante al mezclar $0.5\;L$ de una disolución $2\;M$ de ácido sulfúrico con  $20\;cm^3$ de $\chem{H_2SO_4}$ concentrado al $96\%$ en peso y densidad $1.84\;g\,cm^{-3}$.
		\begin{mdframed}[backgroundcolor=brown!10]  
		Como es la misma sustancia podemos asumir volúmenes aditivos. La disolución resultante tendrá un volumen de $0.52\;L$. De la primera se añaden $0.5\cdot 2\simeq 1\;mol$ de ácido. 
		La masa de disolución procedente de la segunda será $1.84\cdot 20\;g$, que contendrá $0.96\cdot 1.84\cdot 20\;g$ de ácido puro, equivalentes a $\sfrac{0.96\cdot 1.84\cdot  20}{98}\simeq 0.36\;mol$. 
		
		Con todo la disolución resultante tendrá una molaridad
		\[M\simeq \frac{0.36+1}{0.52}\simeq 2.6\;mol\,L^{-1}.\]
		\end{mdframed}
	\end{itemize}
\qed

\newpage
\pagenumbering{gobble}

\question[2]   En condiciones normales de presión y temperatura se hacen reaccionar $60\;g$ de un mineral de cobre de $90\%$ de pureza con $100\;mL$ de una solución $1\,M$ de ácido nítrico. Sabiendo que el rendimiento para la formación de $\chem{NO_2}(g)$ es del $90\%$ y que se forma, además, nitrato de cobre{\small (II)} y agua:
    	\begin{itemize}
    		\item[a)] Ajusta la reacción e identifica el reactivo limitante.
    		\begin{mdframed}[backgroundcolor=brown!10]  
    		La reacción ajustada es
    		\[\chem{Cu+4HNO_3}\longrightarrow\chem{2NO_2+Cu(NO_3)_2+2H_2O}.\]
    		Se requieren, por ello, cuatro veces más moles de ácido nítrico que de cobre. 
    		Calculemos los moles que hay disponibles de cada especie: de cobre tenemos $60\cdot 0.9\simeq 54\;g$ ($0.84\;mol$) y de ácido nítrico, $0.1\cdot 1\simeq 0.1\;mol$. 
    		
    		Naturalmente no están en proporción estequiométrica y va a sobrar cobre. El ácido nítrico es el reactivo limitante.
    		\end{mdframed}
    		\item[b)] Suponiendo que se comporta idealmente, calcula  el volumen de dióxido de nitrógeno que se obtiene.
    	\begin{mdframed}[backgroundcolor=brown!10]  
    	Según la estequiometría de la reacción partir de los $0.1\;mol$ de ácido nítrico se formarán $0.05\;mol$ de dióxido de nitrógeno. Si el rendimiento es del $90\%$ habrán de formarse $0.05\cdot 0.9\simeq 0.045\;mol$, que ocuparán un volumen 
    	\[V=\frac{nRT}{P}=\frac{0.045\cdot 0.082\cdot 273}{1}\simeq 1\;L.\]
    	\end{mdframed}
    	\end{itemize}
  	\qed

\newpage
\pagenumbering{gobble}


\question[2] La descomposición térmica del clorato de potasio se produce según la reacción $\chem{KClO_3(s)}\rightarrow \chem{KCl(s)}+\chem{O_2(g)}.$
	\begin{itemize}
		\item[a)] Calcule la entalpía estándar de reacción para la descomposición de un mol de clorato.
		%-22 kJ/mol de clorato. (ajuste 1-> 1 , 3/2)
		\begin{mdframed}[backgroundcolor=brown!10]  
		La reacción ajustada es 
		\[\chem{KClO_3(s)}\rightarrow \chem{KCl(s)}+\frac{3}{2}\chem{O_2(g)},\]
		de manera que la entalpía de reacción para la descomposición de un mol de clorato es 
		\[\Delta_rH^0(\chem{KClO_3})=\Delta_fH^0(\chem{KCl})-\Delta_fH^0(\chem{KClO_3})\simeq -436-(-414)\simeq -22\;kJ.\]
		\end{mdframed}
		\item[b)] Determine la cantidad de calor, medida a presión constante, que intercambia la reacción con el medio al obtenerse $30\;L$ de oxígeno molecular, medidos en condiciones estándar de presión y temperatura. \textit{¿Se ceden o se absorben?}
		\begin{mdframed}[backgroundcolor=brown!10]  
		La cantidad de calor a presión constante es la entalpía de reacción. Los $30\;L$ de volumen equivalen a $$n=\frac{PV}{RT}=\frac{1\cdot 30}{0.082\cdot 273}\simeq 1.34\;mol,$$
		de manera que según la estequiometría han debido descomponerse $1.34\cdot \sfrac{2}{3}\simeq 0.89\;mol$ de clorato. 
		
		Como por cada mol se desprenden $22\;kJ$, al descomponerse esta cantidad se ceden $0.89\cdot 22\simeq 19.7\;kJ$ al entorno.
		\end{mdframed}
		%se ceden 19.7 kJ al entorno.
	\end{itemize}
\vspace{1ex}
	
	{
	\textbf{Datos:} $\Delta_fH^0(\chem{KClO_3(s)})\simeq -414\;kJ\,mol^{-1},\hspace{2ex}\Delta_fH^0(\chem{KCl(s)})\simeq -436\;kJ\,mol^{-1}.$
	}
\qed
%%
%%%\afterpage{\blankpage}
%%\newpage
%%\vspace*{-7ex}
%%
%%\begin{center}
%%       \fbox{\footnotesize HOJA DE RESPUESTAS A LOS PROBLEMAS}
%%\end{center}

\end{questions}	


\end{document}






